\section{Anexos} \label{sec:anexos}

Los anexos se usan para material complementario que apoya el informe principal sin repetir los ejemplos del documento.
Aquí puedes incluir resultados extendidos, capturas, tablas completas, derivaciones largas o fragmentos de código que sean relevantes para la entrega final.

\subsection{Datos adicionales} \label{sec:anexo_datos_adicionales}
Incluye aquí información que sería demasiado extensa para el cuerpo principal, como:

\begin{itemize}
    \item tablas con resultados completos,
    \item figuras adicionales,
    \item observaciones técnicas o metodológicas,
    \item supuestos o parámetros usados en la implementación.
\end{itemize}

\subsection{Código complementario} \label{sec:anexo_codigo_complementario}
Si necesitas agregar código, procura que sea material que no aparezca ya en la sección de ejemplos.
Por ejemplo, puedes pegar una versión final de un script, una consulta o una función auxiliar que explique mejor el trabajo realizado.

\subsection{Otros recursos} \label{sec:otros_recursos}
Aquí puedes agregar enlaces, referencias internas, resultados de pruebas, o cualquier otro material de respaldo que quieras dejar como soporte del informe.